%%--------------------------------------------------------------
%% template.tex — Example paper using simplejournal class
%%--------------------------------------------------------------
%%
%% CLASS OPTIONS (combine as needed):
%%   onecolumn    — single column (default)
%%   twocolumn    — two-column layout
%%   preprint     — 12pt, double spacing for review
%%   linenumbers  — line numbers in margin
%%
%% Examples:
%%   \documentclass{simplejournal}                         % default: 1-col, 11pt
%%   \documentclass[twocolumn]{simplejournal}              % 2-col, 10pt
%%   \documentclass[preprint,linenumbers]{simplejournal}   % review mode
%%   \documentclass[twocolumn,linenumbers]{simplejournal}  % 2-col + line numbers
%%--------------------------------------------------------------

\documentclass[preprint,linenumbers]{simplejournal}

\title{A Simple Template for Scientific Papers}

\author{
  Alice Author\orcidlink{0000-0002-1825-0097}$^{1}$\email{alice.author@example.com},
  Bob Researcher$^{1,2}$,
  and Carol Scientist$^{2}$\\[6pt]
  \affil{$^1$Department of Physics, University of Examples, Example City}\\
  \affil{$^2$Institute for Advanced Studies, Research Town}
}

\date{\today}

\begin{document}

\maketitle

\begin{abstract}
This document demonstrates the \texttt{simplejournal} class.
It supports one-column and two-column layouts, optional line numbers,
a preprint mode with double spacing, and a clean bibliography style
that sorts references in order of appearance and truncates long author
lists after three names.
No watermarks, no decorations --- just the essentials.
\end{abstract}

\section{Introduction}

This template is intentionally minimal. Modify the class options in the
\verb|\documentclass| line to switch between layouts:
%
\begin{itemize}
  \item \texttt{onecolumn} (default) or \texttt{twocolumn}
  \item \texttt{preprint} for double-spaced review copies
  \item \texttt{linenumbers} to show line numbers
\end{itemize}

We first cite Jain et al.~\cite{Jain2013}, then Wilkinson et al.~\cite{Wilkinson2016},
and finally Kittel~\cite{Randle1997} --- regardless of alphabetical order.

Note that Jain et al.\ has 11 authors but only the first three are printed
in the bibliography, followed by ``et al.'' Wilkinson et al.\ (many authors)
is also truncated the same way.

\section{Stress-test references}

This section cites all entries in the bibliography to verify formatting.
Materials databases and repositories~\cite{Chung2022,Draxl2019,Janssen2019,Huber2020,Merchant2023,Tadmor2011,Hale2016,Walsh2024,BenMahmoud2024,Pyzer-Knapp2022}.
Semantic web and ontologies for materials~\cite{Bayerlein2022,Butler2024,Scheffler2022,Ghiringhelli2023,Valdestilhas2023,Li2020,Bayerlein2024,DelNostro2024,Evans2024,Mrdjenovich2020,Himanen2019,dePablo2019,Horsch2020,Bonomi2019}.
Software infrastructure~\cite{Villamar2025,Rosen2024,Janssen2025,Curtarolo2012,Suarez-Figueroa2012,Menon2019,HjorthLarsen2017,Ong2013,Hale2018,THOMPSON2022,Kresse1996,Giannozzi2009,Hall1991}.
Standards and web resources~\cite{provo,qudt_fairsharing,Poveda2014,wikidata,chebi,owl2_2012,rdf11-concepts,Xu2024,Tudorache2020,Pydantic2026,RDFLib2025,sutton1996interfaces,cmso_git,cmso_zenodo,Patala2019}.

\section{Methods}

Replace this section with your content. Equations work as normal:
\begin{equation}
  E = mc^2
\end{equation}

\section{Results}

Figures, tables, and all standard \LaTeX\ environments are supported.

\section{Conclusions}

Keep it simple.

%% ---- Bibliography ----
\bibliographystyle{simplejournal}
\bibliography{references}

\end{document}
